\section{Stator-flux orientation and direct flux-vector control}
\label{sec:flux-oriented-frames}

The stator-flux frame is an established counterexample to the idea that the
rotor axes must also be the control axes. Let
\(\boldsymbol\psi_s=\psi_s[\cos\delta,\sin\delta]^{\mathsf T}_{dq}\).
Choosing the \(f\) axis along this vector gives
\[
 \boldsymbol\psi_{s,ft}=R(-\delta)\boldsymbol\psi_{s,dq}
 =\begin{bmatrix}\gls{sym:psis}&0\end{bmatrix}^{\mathsf T},\qquad
 \begin{bmatrix}\gls{sym:iflux}\\\gls{sym:itorqueft}\end{bmatrix}
 =R(-\delta)\begin{bmatrix}i_d\\i_q\end{bmatrix}.
\]
The inverse is the transpose rotation. Figure~\ref{fig:arbitrary-frame} makes
the distinction precise: a constant \(\delta\) is merely another rotor-fixed
basis, whereas the \(ft\) angle is a state-dependent flux angle.

\begin{figure}[tb]
\centering
\begin{tikzpicture}[scale=1.05,>=Latex,font=\small]
\coordinate (O) at (0,0);
\draw[->,gray] (O)--(3.1,0) node[right]{\(d\)};
\draw[->,gray] (O)--(0,2.5) node[above]{\(q\)};
\draw[->,MidnightBlue,very thick] (O)--({2.7*cos(32)},{2.7*sin(32)}) node[right]{\(\gamma\)-axis};
\draw[->,MidnightBlue,thick] (O)--({2.0*cos(122)},{2.0*sin(122)}) node[above]{quadrature};
\draw[->,Purple,very thick] (O)--({2.25*cos(60)},{2.25*sin(60)}) node[above]{\(\vect x\)};
\draw[->] (0.75,0) arc[start angle=0,end angle=32,radius=.75];
\node at (0.82,.23) {\(\delta\)};
\node[draw,rounded corners,align=left,fill=Orange!10] at (6.3,1.2)
 {\(\delta=\mathrm{const.}\): rotor-fixed\\
  \(\delta=\delta(\vect i,\boldsymbol\psi)\): moving\\
  extra rate: \(-\dot\delta\matr J\vect x_\gamma\)};
\end{tikzpicture}
\caption{An arbitrary rotor-relative spatial frame. A state-dependent angle
simplifies instantaneous geometry but contributes its own angular velocity.}
\label{fig:arbitrary-frame}
\end{figure}


The air-gap torque is the oriented area spanned by stator flux and current,
\begin{equation}
 M=\gls{sym:kt}(\psi_\alpha i_\beta-\psi_\beta i_\alpha)
 =\boxed{\gls{sym:kt}\gls{sym:psis}\gls{sym:itorqueft}}.
 \label{eq:ft-torque}
\end{equation}
Once one axis lies on the flux vector, the parallel current produces no
oriented area and only the perpendicular component remains. Thus \(i_t\) is
not a newly named \(i_q\): it is a geometrically exact torque-producing
current in the instantaneous stator-flux frame.

\begin{figure}[tb]
\centering
\begin{tikzpicture}[scale=.9,>=Latex,font=\small]
\begin{scope}
 \coordinate (O) at (0,0);
 \draw[->,gray] (O)--(3,0) node[right]{\(d\)};
 \draw[->,gray] (O)--(0,2.6) node[above]{\(q\)};
 \draw[->,ForestGreen!70!black,very thick] (O)--({2.6*cos(38)},{2.6*sin(38)}) node[right]{\(\boldsymbol\psi_s,f\)};
 \draw[->,ForestGreen!70!black,thick] (O)--({1.8*cos(128)},{1.8*sin(128)}) node[above]{\(t\)};
 \draw[->,Purple,very thick] (O)--({2.1*cos(76)},{2.1*sin(76)}) node[above]{\(\vect i_s\)};
 \draw[->] (.65,0) arc[start angle=0,end angle=38,radius=.65];
 \node at (.72,.24){\(\delta\)};
 \draw[dashed,Purple] ({2.1*cos(76)},{2.1*sin(76)})--({1.28*cos(38)},{1.28*sin(38)});
 \node at (2.0,2.35){\(dq\) and \(ft\)};
\end{scope}
\node[draw,rounded corners,align=center,fill=MidnightBlue!7] at (6.3,1.45)
 {\(M_{dq}=k_Ti_q[\psi_{\rm exc}+\Delta L i_d]\)\\[1mm]
  \(\Downarrow\ R(-\delta)\)\\[1mm]
  \(M_{ft}=k_T\psi_s i_t\)};
\node[align=center,text width=4.2cm] at (6.3,-.15)
 {simpler torque and voltage-limit meaning\\at the cost of flux angle estimation};
\end{tikzpicture}
\caption{Rotor \(dq\) and stator-flux \(ft\) frames represent the same spatial
vectors but expose different algebraic relations.}
\label{fig:dq-vs-ft}
\end{figure}

\begin{figure}[tb]
\centering
\begin{tikzpicture}[scale=1.05,>=Latex,font=\small]
\coordinate (O) at (0,0);
\draw[fill=Purple!12,draw=Purple] (O)--(3.3,0)--(2.35,2.05)--cycle;
\draw[->,ForestGreen!70!black,very thick] (O)--(3.3,0) node[right]{\(\boldsymbol\psi_s=\psi_s\boldsymbol e_f\)};
\draw[->,Purple,very thick] (O)--(2.35,2.05) node[above]{\(\vect i=i_f\boldsymbol e_f+i_t\boldsymbol e_t\)};
\draw[dashed,Purple] (2.35,2.05)--(2.35,0);
\draw[<->,Orange,thick] (2.62,0)--(2.62,2.05) node[midway,right]{\(i_t\)};
\node[align=center] at (6.2,1.15)
 {oriented area \(=\psi_s i_t\)\\
  \(i_f\) changes the base projection,\\not the perpendicular height};
\end{tikzpicture}
\caption{Geometric derivation of \(M=(3/2)p\psi_si_t\): torque is the signed
area of the flux--current parallelogram.}
\label{fig:cross-product-torque}
\end{figure}


\subsection{Voltage dynamics in the moving flux frame}

In a frame rotating at
\(\dot\gamma=\omega_e+\dot\delta\), the stator equation is
\[
 \vect u_{ft}=R_s\vect i_{ft}+\dot{\boldsymbol\psi}_{s,ft}
 +(\omega_e+\dot\delta)\matr J\boldsymbol\psi_{s,ft}.
\]
Because the second flux component is zero,
\begin{equation}
 \boxed{u_f=R_si_f+\dot\psi_s},\qquad
 \boxed{u_t=R_si_t+\psi_s(\omega_e+\dot\delta)}.
 \label{eq:ft-voltage}
\end{equation}
The \(f\)-axis voltage changes flux magnitude; \(i_t\) sets torque; and the
price is an observed flux angle and its rotational term. At high speed
\(U\approx\omega_e\psi_s\), so a flux command exposes the voltage limit
directly. A conventional current-vector stack uses
\[
 M^\star\to(i_d^\star,i_q^\star)\to 2\times\mathrm{PI},
\]
whereas flux--torque control uses
\[
 (M^\star,\psi_s^\star)\to(i_t^\star,\psi_s^\star)
 \to\text{torque/current and flux loops}.
\]
The latter may remove a field-weakening current-reference map, but it adds a
stator-flux observer and low-speed magnetic-model dependence.

\subsection{What the literature already establishes}

Unified \gls{dfvc} in stator-flux coordinates was demonstrated across
sinusoidal AC-machine families by Pellegrino, Bojoi, and Guglielmi
\cite{pellegrino2011dfvc}; model-based PMSM control and IPMSM self-learning or
virtual-signal-injection variants further address nonlinear magnetic maps and
MTPA tracking \cite{boazzo2015dfvc,sun2016dfvc,sun2017selflearning}. The EESM
line is no longer a conjectural extension. Ionta et al.\ reported an ECCE
implementation in 2024 \cite{ionta2024dfvc}; the 2026 journal works cover
direct flux/torque control, deadbeat direct-flux control with interchangeable
torque controllers, and a unified design of accurate torque controllers
\cite{ionta2026dftc,rubino2026deadbeat,ionta2026unified}. Metadata checked on
25 September 2026 place the first two journal papers in volume/issue/page
ranges, while the unified torque-controller paper still carries early-access
pages 1--21.

Those works solve practical flux estimation, loop design, torque control, and
experimental realization. The additional question here is different: among
rotor, stator-flux, active-flux, and abstract task coordinates, which chart
best exposes optimization, null-space allocation, conditioning, and embedded
cost? No claim of a new DFVC principle is made.
